% ============================================================
% Chapter 3: Point Estimation — MLE, Method of Moments
% ============================================================
\chapter{Point Estimation --- MLE and Method of Moments}
\label{ch:estimation}

\begin{center}
\textit{``The art of statistics is choosing the right numerical summary.''}
\end{center}

\section{Motivating Example}

Lifetimes (hours) of 10 LED bulbs: $1200, 1350, 980, 1500, 1100, 1450, 1300, 1050, 1400, 1250$. Model as $\mathrm{Exp}(\lambda)$. How to estimate $\lambda$?

\section{General Framework}

\begin{definition}[Estimator]
\index{estimator}
An \textbf{estimator} of $\theta$ is a statistic $\thetahat = T(X_1,\ldots,X_n)$ taking values in $\Theta$.
\end{definition}

\section{Method of Moments (MoM)}
\index{method of moments}

\begin{definition}[Method of Moments]
Equate the first $k=\dim(\theta)$ population moments to sample moments:
\[
\E_\theta[X^j] = \frac{1}{n}\sum_{i=1}^n X_i^j, \quad j=1,\ldots,k.
\]
\end{definition}

\begin{example}[Normal distribution]
$\hat{\mu}_{\mathrm{MoM}} = \xbar$, $\hat{\sigma}^2_{\mathrm{MoM}} = \frac{1}{n}\sum(X_i-\xbar)^2$.
\end{example}

\begin{example}[Exponential distribution]
$\E[X]=1/\lambda$, so $\hat{\lambda}_{\mathrm{MoM}} = 1/\xbar$.
\end{example}

\section{Maximum Likelihood Estimation (MLE)}
\index{MLE}\index{maximum likelihood}

\begin{definition}[Likelihood Function]
\[
L(\theta) = \prod_{i=1}^n f(x_i;\theta), \qquad \ell(\theta)=\ln L(\theta) = \sum_{i=1}^n \ln f(x_i;\theta).
\]
\end{definition}

\begin{definition}[Maximum Likelihood Estimator]
\[
\thetahat_{\mathrm{MLE}} = \argmax_{\theta\in\Theta} \ell(\theta).
\]
\end{definition}

\subsection{Fundamental Examples}

\begin{example}[Bernoulli]
$\phat_{\mathrm{MLE}} = \xbar$.
\end{example}

\begin{example}[Normal]
$\hat{\mu}_{\mathrm{MLE}} = \xbar$, $\hat{\sigma}^2_{\mathrm{MLE}} = \frac{1}{n}\sum(X_i-\xbar)^2$.
\end{example}

\begin{warning}
The MLE of $\sigma^2$ uses divisor $n$, which is \textbf{biased}. The corrected version $S^2$ (divisor $n-1$) is unbiased (cf.\ Chapter~\ref{ch:properties}).
\end{warning}

\begin{example}[Exponential]
$\hat{\lambda}_{\mathrm{MLE}}=1/\xbar$.
\end{example}

\begin{example}[Poisson]
$\hat{\lambda}_{\mathrm{MLE}}=\xbar$.
\end{example}

\begin{example}[Uniform {$\mathrm{Unif}([0,\theta])$}]
$L(\theta)=\theta^{-n}\mathbf{1}_{\theta\geq x_{(n)}}$, decreasing for $\theta\geq x_{(n)}$. So $\thetahat_{\mathrm{MLE}}=X_{(n)}=\max_i X_i$. The score equation does not apply (boundary maximum).
\end{example}

\section{Properties of the MLE}

\begin{theorem}[Invariance]
\index{invariance}
If $\thetahat$ is the MLE of $\theta$, then $g(\thetahat)$ is the MLE of $g(\theta)$.
\end{theorem}

\begin{theorem}[Consistency]
Under regularity conditions, $\thetahat_{\mathrm{MLE}} \xrightarrow{\Prob} \theta_0$.
\end{theorem}

\begin{theorem}[Asymptotic Normality]
\label{thm:mle-normality}
Under regularity conditions:
\[
\sqrt{n}(\thetahat_{\mathrm{MLE}}-\theta_0) \xrightarrow{\mathcal{L}} \mathcal{N}\!\left(0, \frac{1}{I(\theta_0)}\right).
\]
\end{theorem}

\section{Fisher Information}
\index{Fisher information}

\begin{definition}[Fisher Information]
\[
I(\theta) = \E_\theta\!\left[\left(\frac{\partial \ln f(X;\theta)}{\partial\theta}\right)^{\!2}\right] = -\E_\theta\!\left[\frac{\partial^2 \ln f(X;\theta)}{\partial\theta^2}\right].
\]
For $n$ observations: $I_n(\theta)=nI(\theta)$.
\end{definition}

\begin{example}
$X\sim\mathcal{N}(\mu,\sigma^2)$ ($\sigma^2$ known): $I(\mu)=1/\sigma^2$.
\end{example}

\section{Python Implementation}

\begin{python}
\begin{minted}{python}
import numpy as np
from scipy import stats
import matplotlib.pyplot as plt

data = np.array([1200, 1350, 980, 1500, 1100, 1450, 1300, 1050, 1400, 1250])

# MoM and MLE for Exp(lambda)
lambda_hat = 1 / np.mean(data)
print(f"lambda_hat = {lambda_hat:.6f}")

# MLE for Normal
mu_mle = np.mean(data)
sigma2_mle = np.var(data, ddof=0)
sigma2_unbiased = np.var(data, ddof=1)
print(f"Normal MLE: mu={mu_mle:.2f}, sigma^2={sigma2_mle:.2f}")
print(f"Unbiased:   sigma^2={sigma2_unbiased:.2f}")

# Log-likelihood plot
lambdas = np.linspace(0.0005, 0.002, 200)
log_lik = [np.sum(stats.expon.logpdf(data, scale=1/lam)) for lam in lambdas]
plt.figure(figsize=(8, 5))
plt.plot(lambdas, log_lik, 'b-', lw=2)
plt.axvline(lambda_hat, color='r', ls='--', label=f'MLE = {lambda_hat:.5f}')
plt.xlabel(r'$\lambda$'); plt.ylabel(r'$\ell(\lambda)$')
plt.title('Exponential log-likelihood'); plt.legend()
plt.savefig("ch03_loglik.pdf"); plt.show()

# Fisher information
I_fisher = 1 / lambda_hat**2
var_asymp = 1 / (len(data) * I_fisher)
print(f"Asymptotic SE: {np.sqrt(var_asymp):.6f}")
\end{minted}
\end{python}

\section{Exercises}

\begin{exercise}[$\star$]
Compute the MLE for (a) Geometric$(p)$, (b) $\mathrm{Unif}([a,b])$ with both $a,b$ unknown.
\end{exercise}

\begin{exercise}[$\star$]
Find MoM estimators for $\mathrm{Beta}(\alpha,\beta)$.
\end{exercise}

\begin{exercise}[$\star\star$ -- Invariance]
Under $\mathcal{N}(\mu,\sigma^2)$, find the MLE of $\Prob(X>c)$.
\end{exercise}

\begin{exercise}[$\star\star$ -- Fisher information]
Compute $I(\theta)$ for Bernoulli, Poisson, and the Fisher information matrix for $\mathcal{N}(\mu,\sigma^2)$.
\end{exercise}

\begin{exercise}[$\star\star\star$ -- Project: MoM vs MLE comparison]
For $\Gamma(\alpha,\beta)$: simulate and compare bias and MSE of MoM vs MLE for $n=20,50,200$.
\end{exercise}

\begin{keyformulas}
\begin{align*}
L(\theta) &= \prod f(x_i;\theta), \quad \ell(\theta)=\sum \ln f(x_i;\theta) \\
\thetahat_{\mathrm{MLE}} &= \argmax_\theta \ell(\theta) \\
I(\theta) &= -\E\!\left[\frac{\partial^2\ln f}{\partial\theta^2}\right] \\
\sqrt{n}(\thetahat_{\mathrm{MLE}}-\theta) &\xrightarrow{\mathcal{L}} \mathcal{N}(0,1/I(\theta))
\end{align*}
\end{keyformulas}
